\documentclass[a4paper,12pt]{article}
	\usepackage[ansinew]{inputenc}
	\usepackage{ae}
	\usepackage[brazil]{babel}
	\usepackage{lipsum}
	\usepackage{xcolor}	
	\usepackage{graphicx}	
	\usepackage{subfigure}
	
	\graphicspath{
		{Figuras/EPS/}
		{Figuras/JPG/}
		{Figuras/PDF/}
		{Figuras/PNG/}
		}
		
	\newcommand{\comando}[1]{{\normalfont\textbackslash\color{blue!50!black}#1}}
\newcommand{\pacote}[1]{{\normalfont\sffamily\color{orange!50!black}#1}}
\newcommand{\ambiente}[1]{{\normalfont\sffamily\color{green!50!black}#1}}
\newcommand{\contador}[1]{{\normalfont\sffamily\color{red!50!black}#1}}
\newcommand{\foreign}[1]{\textsl{#1}}
\newcommand\linhabase{\noindent\makebox[0pt][l]{\color{red}\rule{\textwidth}{0.1pt}}}
	
\begin{document}

	\setcounter{lofdepth}{2}
	\listoffigures
	
	\newpage

	\lipsum[1]

	\begin{figure}[htb]
		\center
		\subfigure
			[Esta é a legenda que aparece na lista de figuras (LoF)]
			[\label{fig:figA}Esta é a legenda longa da figura \ref{fig:figA}. Se você fornecer uma legenda curta no argumento opcional anterior, esta legenda \emph{não} aparecerá na lista de figuras (LoF), mas apenas no texto\dots]
			{\includegraphics[width=0.45\textwidth]{Tei-Gi}}
		\qquad
		\subfigure
			[Esta é a legenda da figura \ref{fig:figB} que aparece na LoF.]
			[\dots\ Ainda assim é necessário alterar o valor do contador \contador{lofdepth} para $2$ de modo que as legendas das sub-figuras apareçam na LoF. Caso contrário, apenas a legenda do conjunto (no ambiente \ambiente{figure}) aparecerá. Este contador, \foreign{list of figure depth}, marca a ``profundidade'' da LoF: $\leq1$ são figuras e $\geq2$ são sub-figuras.]
			{\label{fig:figB}\includegraphics[width=0.45\textwidth]{Tei-Gi}}
			
		\caption
			[Esta é a legenda comum às sub-figuras que aparece na LoF]
			{observe que no comando \comando{subfigure} você pode colocar o \comando{label} tanto no argumento opcional da legenda (a que aparece no texto) como no argumento obrigatório, que contém o comando \comando{includegraphics}.}
	\end{figure}

	\lipsum[2]

\end{document}